\documentclass[10pt,a4paper]{amsart}
\usepackage{amsmath,amssymb,amsthm,graphicx}
% Retained preamble for the Stacks Project chapter snapshot in this directory.
% The chapter's own preamble.tex is a copy of the shared Stacks Project preamble; the block
% below is sliced MECHANICALLY, line for line, from it. Nothing was added, removed or
% redefined; the wrapper class, the xr-hyper/hyperref wiring and the externaldocument list
% are omitted because the assembler supplies the document class.
%
% lines taken from preamble.tex: 3, 5, 9, 11, 12, 156, 157, 158, ... (51 lines total)
%
% package coverage restored from the chapter's own preamble (2026-10-13): xy, xy-2cell

\usepackage[all]{xy}
\xyoption{2cell}
\UseAllTwocells

\newenvironment{reference}{\comment}{\endcomment}
\newenvironment{slogan}{\comment}{\endcomment}
\newenvironment{history}{\comment}{\endcomment}
\theoremstyle{plain}
\newtheorem{theorem}[subsection]{Theorem}
\newtheorem{proposition}[subsection]{Proposition}
\newtheorem{lemma}[subsection]{Lemma}
\theoremstyle{definition}
\newtheorem{definition}[subsection]{Definition}
\newtheorem{example}[subsection]{Example}
\newtheorem{exercise}[subsection]{Exercise}
\newtheorem{situation}[subsection]{Situation}
\theoremstyle{remark}
\newtheorem{remark}[subsection]{Remark}
\newtheorem{remarks}[subsection]{Remarks}
\numberwithin{equation}{subsection}
\def\lim{\mathop{\mathrm{lim}}\nolimits}
\def\colim{\mathop{\mathrm{colim}}\nolimits}
\def\Spec{\mathop{\mathrm{Spec}}}
\def\Hom{\mathop{\mathrm{Hom}}\nolimits}
\def\Ext{\mathop{\mathrm{Ext}}\nolimits}
\def\SheafHom{\mathop{\mathcal{H}\!\mathit{om}}\nolimits}
\def\SheafExt{\mathop{\mathcal{E}\!\mathit{xt}}\nolimits}
\def\Sch{\mathit{Sch}}
\def\Mor{\mathop{\mathrm{Mor}}\nolimits}
\def\Ob{\mathop{\mathrm{Ob}}\nolimits}
\def\Sh{\mathop{\mathit{Sh}}\nolimits}
\def\NL{\mathop{N\!L}\nolimits}
\def\CH{\mathop{\mathrm{CH}}\nolimits}
\def\proetale{{pro\text{-}\acute{e}tale}}
\def\etale{{\acute{e}tale}}
\def\QCoh{\mathit{QCoh}}
\def\Ker{\mathop{\mathrm{Ker}}}
\def\Im{\mathop{\mathrm{Im}}}
\def\Coker{\mathop{\mathrm{Coker}}}
\def\Coim{\mathop{\mathrm{Coim}}}
\DeclareMathSymbol{\boxtimes}{\mathbin}{AMSa}{"02}
\def\QCohstack{\mathcal{QC}\!\mathit{oh}}
\def\Cohstack{\mathcal{C}\!\mathit{oh}}
\def\Spacesstack{\mathcal{S}\!\mathit{paces}}
\def\Quotfunctor{\mathrm{Quot}}
\def\Hilbfunctor{\mathrm{Hilb}}
\def\Curvesstack{\mathcal{C}\!\mathit{urves}}
\def\Polarizedstack{\mathcal{P}\!\mathit{olarized}}
\def\Complexesstack{\mathcal{C}\!\mathit{omplexes}}
\def\Pic{\mathop{\mathrm{Pic}}\nolimits}
\def\Picardstack{\mathcal{P}\!\mathit{ic}}
\def\Picardfunctor{\mathrm{Pic}}
\def\Deformationcategory{\mathcal{D}\!\mathit{ef}}
\usepackage{cleveref}
\crefname{originallemma}{original lemma}{original lemmas}
\Crefname{originallemma}{Original lemma}{Original lemmas}
\makeatletter
\newlabel{simplicial-example-fibre-products-simplicial-object}{{016E}{source}{}{sourcepointer-1}{}}
\newlabel{simplicial-example-fibre-products-simplicial-object@cref}{{[originallemma][0][]016E}{[1][0][]source}{}{}{}}
\makeatother
\crefname{originallemma}{original lemma}{original lemmas}
\Crefname{originallemma}{Original lemma}{Original lemmas}
\makeatletter
\newlabel{simplicial-lemma-cosk-minus-one-equivalence}{{01AB}{source}{}{sourcepointer-2}{}}
\newlabel{simplicial-lemma-cosk-minus-one-equivalence@cref}{{[originallemma][0][]01AB}{[1][0][]source}{}{}{}}
\makeatother
\crefname{originallemma}{original lemma}{original lemmas}
\Crefname{originallemma}{Original lemma}{Original lemmas}
\makeatletter
\newlabel{simplicial-lemma-homology-eilenberg-maclane}{{0197}{source}{}{sourcepointer-3}{}}
\newlabel{simplicial-lemma-homology-eilenberg-maclane@cref}{{[originallemma][0][]0197}{[1][0][]source}{}{}{}}
\makeatother
\crefname{originallemma}{original lemma}{original lemmas}
\Crefname{originallemma}{Original lemma}{Original lemmas}
\makeatletter
\newlabel{simplicial-lemma-trivial-kan-homotopy}{{08NS}{source}{}{sourcepointer-4}{}}
\newlabel{simplicial-lemma-trivial-kan-homotopy@cref}{{[originallemma][0][]08NS}{[1][0][]source}{}{}{}}
\makeatother
\title{Homology of the Cech nerve of a sheaf-surjective map of presheaves}
\author{The Stacks Project Authors}
\date{stacks-project:hypercovering.tex:local-snapshot-20261012; GNU FDL 1.2 or later}
\begin{document}
\maketitle
\noindent\textit{Editorial note.} Statement span is lines 807-831 of sources/stacks-hypercovering/pinned-original-source.tex; proof span is lines 833-877, the complete author proof (43 body lines).\par\medskip
\noindent\textit{Editorial note.} Navigation stubs: 4 cited-but-not-included labels are supplied as non-load-bearing pointers with their pinned tags where the snapshot tag index has them.\par\medskip
\noindent\textit{Editorial note.} extra\_packages is empty on purpose: the retained preamble already loads xy with options and a second declaration clashes in the assembler wrapper.\par\medskip
\noindent\textit{Editorial note.} Text is retained verbatim, including the author's own typos and placeholders.\par\medskip
\noindent\textit{Editorial note.} Queue index 11 of queue\_hypercovering.json (pb=1673, stmt 807-831, proof 833-877).\par\medskip
\noindent\textit{Editorial note.} stacks-hypercovering wave-1 (hc01): fresh adjudication accept (2026-10-14); pool per standing fallback order (hypercovering -> pic -> spaces-more-cohomology) after md01 closed the moduli pool; hypercovering.tex 53 labels/34 with proof judged 26 accept / 8 reject / 14 pb=0 env-dup rows excluded; spans cite-free and reference/slogan/history-free so no patchenv needed.\par\medskip
\noindent\textit{Allowed standard prerequisites (editorial).}\begin{itemize}
\item Explicit prerequisite boundary. The cited inputs are prerequisites; the argument in the retained span is not replaced by any cited result.
\end{itemize}\medskip
\ifdefined\hypertarget\hypertarget{sourcepointer-1}{}\fi
\noindent\textit{Source reference.} Original lemma 016E: sources/stacks-hypercovering/pinned-original-source.tex, source lines 877--877. Cited by the retained statement or proof but not part of the collected spans; supplied only as a navigation pointer. The argument does not depend on its content.\par\medskip
\ifdefined\hypertarget\hypertarget{sourcepointer-2}{}\fi
\noindent\textit{Source reference.} Original lemma 01AB: sources/stacks-hypercovering/pinned-original-source.tex, source lines 877--877. Cited by the retained statement or proof but not part of the collected spans; supplied only as a navigation pointer. The argument does not depend on its content.\par\medskip
\ifdefined\hypertarget\hypertarget{sourcepointer-3}{}\fi
\noindent\textit{Source reference.} Original lemma 0197: sources/stacks-hypercovering/pinned-original-source.tex, source lines 877--877. Cited by the retained statement or proof but not part of the collected spans; supplied only as a navigation pointer. The argument does not depend on its content.\par\medskip
\ifdefined\hypertarget\hypertarget{sourcepointer-4}{}\fi
\noindent\textit{Source reference.} Original lemma 08NS: sources/stacks-hypercovering/pinned-original-source.tex, source lines 877--877. Cited by the retained statement or proof but not part of the collected spans; supplied only as a navigation pointer. The argument does not depend on its content.\par\medskip
\begin{lemma}
\label{lemma-compare-cosk0}
Let $\mathcal{C}$ be a site.
Let $\mathcal{F} \to \mathcal{G}$ be a morphism
of presheaves of sets. Denote $K$ the simplicial
object of $\textit{PSh}(\mathcal{C})$ whose $n$th
term is the $(n + 1)$st fibre product of $\mathcal{F}$
over $\mathcal{G}$, see
Simplicial, Example \ref{simplicial-example-fibre-products-simplicial-object}.
Then, if $\mathcal{F} \to \mathcal{G}$ is surjective after
sheafification, we have
$$
H_i(K) =
\left\{
\begin{matrix}
0 & \text{if} & i > 0\\
\mathbf{Z}_\mathcal{G}^\# & \text{if} & i = 0
\end{matrix}
\right.
$$
The isomorphism in degree $0$ is given by the
morphism $H_0(K) \to \mathbf{Z}_\mathcal{G}^\#$
coming from the map $(\mathbf{Z}_K^\#)_0 =
\mathbf{Z}_\mathcal{F}^\# \to \mathbf{Z}_\mathcal{G}^\#$.
\end{lemma}

\begin{proof}
Let $\mathcal{G}' \subset \mathcal{G}$ be the image of
the morphism $\mathcal{F} \to \mathcal{G}$.
Let $U \in \Ob(\mathcal{C})$. Set
$A = \mathcal{F}(U)$ and $B = \mathcal{G}'(U)$.
Then the simplicial set $K(U)$ is equal to the simplicial
set with $n$-simplices given by
$$
A \times_B A \times_B \ldots \times_B A\ (n + 1 \text{ factors)}.
$$
By Simplicial, Lemma \ref{simplicial-lemma-cosk-minus-one-equivalence}
the morphism $K(U) \to B$ is a trivial Kan fibration.
Thus it is a homotopy equivalence
(Simplicial, Lemma \ref{simplicial-lemma-trivial-kan-homotopy}).
Hence applying the functor ``free abelian group on'' to this
we deduce that
$$
\mathbf{Z}_K(U) \longrightarrow \mathbf{Z}_B
$$
is a homotopy equivalence. Note that $s(\mathbf{Z}_B)$ is
the complex
$$
\ldots \to
\bigoplus\nolimits_{b \in B}\mathbf{Z} \xrightarrow{0}
\bigoplus\nolimits_{b \in B}\mathbf{Z} \xrightarrow{1}
\bigoplus\nolimits_{b \in B}\mathbf{Z} \xrightarrow{0}
\bigoplus\nolimits_{b \in B}\mathbf{Z} \to 0
$$
see Simplicial, Lemma \ref{simplicial-lemma-homology-eilenberg-maclane}.
Thus we see that
$H_i(s(\mathbf{Z}_K(U))) = 0$ for $i > 0$, and
$H_0(s(\mathbf{Z}_K(U))) = \bigoplus_{b \in B}\mathbf{Z}
= \bigoplus_{s \in \mathcal{G}'(U)} \mathbf{Z}$.
These identifications are compatible with restriction
maps.

\medskip\noindent
We conclude that $H_i(s(\mathbf{Z}_K)) = 0$ for $i > 0$ and
$H_0(s(\mathbf{Z}_K)) = \mathbf{Z}_{\mathcal{G}'}$, where here
we compute homology groups in $\textit{PAb}(\mathcal{C})$. Since
sheafification is an exact functor we deduce the result
of the lemma. Namely, the exactness implies
that $H_0(s(\mathbf{Z}_K))^\# = H_0(s(\mathbf{Z}_K^\#))$,
and similarly for other indices.
\end{proof}

\end{document}
